\section{Task 2: Feed Forward Neural Network}

\begin{questionbox}
\textbf{Q1.} Do you have the same number of API calls for each sequence? If not, is the distribution
of the number of API calls per sequence in the training set the same as the test one?
\end{questionbox}

No. The API sequences do not all have the same length. In the training set, sequence lengths range
from 60 to 90 API calls, with mean 75.031 and median 75. In the test set, lengths range from 70 to
100, with mean 86.331 and median 87. Therefore, the test distribution is shifted toward longer
sequences.

\begin{table}[h!]
\centering
\begin{tabular}{lcccccccc}
\hline
\textbf{Split} & \textbf{Samples} & \textbf{Min} & \textbf{P25} & \textbf{Median} & \textbf{Mean} & \textbf{P75} & \textbf{P95} & \textbf{Max} \\
\hline
Train & 16325 & 60 & 67.0 & 75.0 & 75.031 & 83.0 & 89.0 & 90 \\
Test  & 6505  & 70 & 79.0 & 87.0 & 86.331 & 94.0 & 99.0 & 100 \\
\hline
\end{tabular}
\caption{Sequence length statistics for train and test partitions.}
\label{tab:task2_sequence_lengths}
\end{table}

\begin{figure}[h!]
\centering
\includegraphics[width=0.70\textwidth]{figures/task-2/sequence_lengths.pdf}
\caption{Distribution of API sequence lengths in train and test.}
\label{fig:task2_sequence_lengths}
\end{figure}

\begin{questionbox}
\textbf{Q2.} Can a FFNN handle a variable number of elements? If not, why?
\end{questionbox}

No. A standard FFNN expects each sample to be represented by a fixed-size input vector. Since the
first linear layer has a fixed number of input units, all sequences must be converted to the same
length before being passed to the network.

\begin{questionbox}
\textbf{Q3.} How to estimate a fixed-size candidate? What partition do you use to estimate it?
\end{questionbox}

The fixed length was estimated using only the training partition, to avoid using test information in
preprocessing. We chose the maximum training sequence length, which is 90 API calls. This keeps all
training sequences intact.

\begin{questionbox}
\textbf{Q4.} Given the estimate of the previous point, what technique could you use to obtain the
same number of API calls per sequence?
\end{questionbox}

We used padding and truncation. Sequences shorter than 90 API calls are padded with the special token
\texttt{<PAD>} mapped to ID 0. API calls not present in the training vocabulary are mapped to
\texttt{<UNK>}. The vocabulary contains 260 IDs in total: 258 API calls plus \texttt{<PAD>} and
\texttt{<UNK>}.

\begin{questionbox}
\textbf{Q5.} Suppose that at test time you have more API calls than the fixed-size you estimated.
What do you do with the API calls exceeding such fixed-size?
\end{questionbox}

Test sequences longer than 90 API calls are truncated. In the test set, 2,639 sequences were longer
than the fixed size, for a total of 14,601 truncated API calls. There were also 3 unknown API calls,
corresponding to the API calls that appear only in the test vocabulary.

\begin{table}[h!]
\centering
\begin{tabular}{lccccc}
\hline
\textbf{Split} & \textbf{Samples} & \textbf{Fixed length} & \textbf{Vocab size} &
\textbf{Truncated seq.} & \textbf{Unknown calls} \\
\hline
Train & 16325 & 90 & 260 & 0    & 0 \\
Test  & 6505  & 90 & 260 & 2639 & 3 \\
\hline
\end{tabular}
\caption{Padding, truncation, and unknown-token statistics.}
\label{tab:task2_preprocessing}
\end{table}

The API calls are categorical strings, so we tested two numerical representations:

\begin{itemize}
    \item \textbf{Sequential identifiers}: each API call is mapped to a fixed arbitrary ID. The padded
    ID vector is normalized and used directly as the FFNN input.
    \item \textbf{Learnable embeddings}: each API ID is mapped to a dense vector learned during
    training; the concatenated embeddings are passed to the FFNN.
\end{itemize}

\begin{questionbox}
\textbf{Q6.} Use a FFNN in both cases. Report how you selected the hyperparameters of your final
model, and justify your choices.
\end{questionbox}

For both FFNN variants, we used Optuna on a stratified validation split from the training set. The
objective was macro F1, because the dataset is strongly imbalanced toward malware. The search space
is reported in Table~\ref{tab:task2_optuna_space}. Both models used weighted cross-entropy to reduce
the effect of class imbalance and early stopping based on validation macro F1.

\begin{table}[h!]
\centering
\begin{tabular}{llll}
\hline
\textbf{Hyperparameter} & \textbf{Search space} & \textbf{Scale/type} & \textbf{Models} \\
\hline
\texttt{hidden\_dim} & \{64, 128, 256\} & Categorical & Both \\
\texttt{dropout} & $[0, 0.5]$ & Uniform & Both \\
\texttt{batch\_size} & \{128, 256, 512\} & Categorical & Both \\
\texttt{lr} & $[10^{-4}, 3 \cdot 10^{-3}]$ & Log-uniform & Both \\
\texttt{weight\_decay} & $[10^{-6}, 10^{-2}]$ & Log-uniform & Both \\
\texttt{embedding\_dim} & \{8, 16, 32\} & Categorical & Embedding only \\
\hline
\end{tabular}
\caption{Optuna search space for the FFNN models.}
\label{tab:task2_optuna_space}
\end{table}

\begin{table}[h!]
\centering
\begin{tabular}{lcccccc}
\hline
\textbf{Model} & \textbf{Hidden} & \textbf{Dropout} & \textbf{LR} & \textbf{Weight decay} &
\textbf{Batch} & \textbf{Emb. dim} \\
\hline
Sequential-ID FFNN & 256 & 0.0697 & 0.000270 & 0.000029 & 256 & -- \\
Embedding FFNN     & 256 & 0.0806 & 0.002362 & 0.001708 & 256 & 16 \\
\hline
\end{tabular}
\caption{Best hyperparameters selected with Optuna for the two FFNN variants.}
\label{tab:task2_hyperparameters}
\end{table}

\begin{figure}[h!]
\centering
\begin{minipage}{0.48\textwidth}
\centering
\includegraphics[width=\textwidth]{figures/task-2/ids_training_curve.pdf}
\end{minipage}
\hfill
\begin{minipage}{0.48\textwidth}
\centering
\includegraphics[width=\textwidth]{figures/task-2/embedding_training_curve.pdf}
\end{minipage}
\caption{Validation macro F1 curves for the Sequential-ID FFNN (left) and the Embedding FFNN (right).}
\label{fig:task2_training_curves}
\end{figure}

The best validation macro F1 was 0.5403 for the Sequential-ID FFNN and 0.6838 for the Embedding
FFNN. On the test set, the embedding model also performed better overall.

\begin{table}[h!]
\centering
\begin{tabular}{lccccc}
\hline
\textbf{Model} & \textbf{Accuracy} & \textbf{Balanced acc.} & \textbf{Macro F1} &
\textbf{F1 Goodware} & \textbf{F1 Malware} \\
\hline
Sequential-ID FFNN & 0.9397 & 0.5138 & 0.5155 & 0.0622 & 0.9689 \\
Embedding FFNN     & 0.9356 & 0.5789 & 0.5749 & 0.1832 & 0.9665 \\
\hline
\end{tabular}
\caption{Final test performance of the two FFNN variants.}
\label{tab:task2_test_performance}
\end{table}

\begin{table}[h!]
\centering
\begin{tabular}{lcccc}
\hline
\textbf{Model/Class} & \textbf{Precision} & \textbf{Recall} & \textbf{F1-score} & \textbf{Support} \\
\hline
Seq-ID Goodware (0) & 0.0743 & 0.0535 & 0.0622 & 243 \\
Seq-ID Malware (1)  & 0.9637 & 0.9741 & 0.9689 & 6262 \\
Emb. Goodware (0)   & 0.1741 & 0.1934 & 0.1832 & 243 \\
Emb. Malware (1)    & 0.9686 & 0.9644 & 0.9665 & 6262 \\
\hline
\end{tabular}
\caption{Per-class classification report for the two FFNN variants.}
\label{tab:task2_classification_report}
\end{table}

\begin{figure}[h!]
\centering
\begin{minipage}{0.48\textwidth}
\centering
\includegraphics[width=\textwidth]{figures/task-2/ids_confusion_matrix.pdf}
\end{minipage}
\hfill
\begin{minipage}{0.48\textwidth}
\centering
\includegraphics[width=\textwidth]{figures/task-2/embedding_confusion_matrix.pdf}
\end{minipage}
\caption{Confusion matrices for the Sequential-ID FFNN (left) and the Embedding FFNN (right).}
\label{fig:task2_confusion_matrices}
\end{figure}

\begin{questionbox}
\textbf{Q7.} Can you obtain the same results for the two alternatives (sequential identifiers and
learnable embeddings)? If not, why?
\end{questionbox}

No. The embedding-based FFNN achieved better macro F1, balanced accuracy, and goodware F1.
Sequential IDs are arbitrary categorical labels, but when used directly as numeric inputs the network
can interpret their magnitude as meaningful even though ID order has no semantic meaning. Learnable
embeddings are more appropriate because the model can learn a dense representation for each API call.
Even so, both FFNNs still struggle on the minority goodware class and both are worse than the
frequency-based Logistic Regression baseline in macro F1.
